\documentclass[a4paper]{article}

\usepackage[utf8x]{inputenc}    
\usepackage[T1]{fontenc}
\usepackage[francais,bloc]{automultiplechoice}    
% [nowatermark] pour enlever le texte "PROJET" et le pied de page
\usepackage{multicol}
\usepackage{fancyhdr,amssymb,amsmath}
\usepackage{pythontex}
%\usepackage{fancyvrb, fvextra, etoolbox, xstring, pgfopts, newfloat, currfile, color}
%===================
\begin{document}
%===================

%%% On répartie les quesitons dans différents groupes.
% ---------------------------------
%%% Questions groupe : "catégorie1"
%----------------------------------
\element{categorie1}{
  \begin{question}{calcul littéral 1}    
% Question ouverte   
Développer et réduire l'expression suivante $(2x-1)^2$  
   
   \AMCOpen{lines=2,dots=false}{\wrongchoice{NA}\scoring{0}\wrongchoice{AB}\scoring{1}\wrongchoice{B}\scoring{2}\correctchoice{TB}\scoring{3}}
% AB rapporte 1 point
% B rapporte 2 points
% TB rapporte 3 points     
  \end{question}
}

%
\element{categorie1}{
  \begin{question}{calcul littéral 2}  
% Question ouverte  
Développer et réduire l'expression suivante $(2x-4)(-x+3)$  
   \AMCOpen{lines=2,dots=false}{\wrongchoice{NA}\scoring{0}\wrongchoice{AB}\scoring{1}\wrongchoice{B}\scoring{2}\correctchoice{TB}\scoring{3}}
% AB rapporte 1 point
% B rapporte 2 points
% TB rapporte 3 points         
  \end{question}
}





% ---------------------------------
%%% Questions groupe : "catégorie2"
%----------------------------------
\element{categorie2}{
  \begin{question}{calcul littéral 3}  
% Question ouverte  
Factoriser l'expression suivante  $(2x-2)(4+5x)+(2x-2)(-x+7)$
   \AMCOpen{lines=2,dots=false}{\wrongchoice{NA}\scoring{0}\wrongchoice{AB}\scoring{1}\wrongchoice{B}\scoring{2}\correctchoice{TB}\scoring{3}}
% AB rapporte 1 point
% B rapporte 2 points
% TB rapporte 3 points         
  \end{question}
}

\element{categorie2}{
  \begin{question}{calcul littéral 4}  
% Question ouverte  
Factoriser l'expression suivante  $(x-2)(4+6x)+(x-2)(-5x+7)$
   \AMCOpen{lines=2,dots=false}{\wrongchoice{NA}\scoring{0}\wrongchoice{AB}\scoring{1}\wrongchoice{B}\scoring{2}\correctchoice{TB}\scoring{3}}
% AB rapporte 1 point
% B rapporte 2 points
% TB rapporte 3 points         
  \end{question}
}
% ---------------------------------
%%% Questions groupe : "catégorie3"
%----------------------------------





\element{categorie3}{
  \begin{question}{calcul terme ari }  
% Question simple à réponse unique    
Soit $(V_n)$ une suite arithmétique telle que $V_5=2$ et de raison $6$, déterminer $V_7$.  

     \begin{reponseshoriz}\bareme{b=1,m=-0.5,e=-0.5,v=0}
      \bonne{$V_7=14$}
      \mauvaise{$V_7=72$}
      \mauvaise{$V_7=-10$}
      \mauvaise{$V_7=\dfrac{1}{18}$}
    \end{reponseshoriz} 
  \end{question}
% b = 1 : une bonne réponse rapporte 1 point
% m = -0.5 : une mauvaise réponse retire 0.5 point
% e = -0.5 : une réponse incohérente retire 0.5 point
% v = 0 : 0 point si aucune réponse  
}


\element{categorie3}{
  \begin{question}{Calcul terme geo}  
% Question simple à réponse unique    
Soit $(U_n)$ une suite géométrique telle que $U_5=2$ et de raison $6$, déterminer $U_7$. 

     \begin{reponseshoriz}\bareme{b=1,m=-0.5,e=-0.5,v=0}
      \mauvaise{$U_7=52$}
      \bonne{$U_7=72$}
      \mauvaise{$U_7=\dfrac{1}{3}$}
      \mauvaise{$U_7=12$}
    \end{reponseshoriz} 
  \end{question}
% b = 1 : une bonne réponse rapporte 1 point
% m = -0.5 : une mauvaise réponse retire 0.5 point
% e = -0.5 : une réponse incohérente retire 0.5 point
% v = 0 : 0 point si aucune réponse  
}

\element{categorie3}{
  \begin{question}{variation ari}  
% Question simple à réponse unique    
Soit $(W_n)$ une suite arithmétique telle que $W_0=-2$ et de raison $6$, on dit que les variations de $(W_n)$ sont sur $\mathbb{N}$.

     \begin{reponseshoriz}\bareme{b=1,m=-0.5,e=-0.5,v=0}
      \mauvaise{décroissante}
      \bonne{strictement croissante}
      \mauvaise{strictement décroissante}
      \mauvaise{constante}
    \end{reponseshoriz} 
  \end{question}
% b = 1 : une bonne réponse rapporte 1 point
% m = -0.5 : une mauvaise réponse retire 0.5 point
% e = -0.5 : une réponse incohérente retire 0.5 point
% v = 0 : 0 point si aucune réponse  
}


\element{categorie3}{
  \begin{question}{Variation geo}  
% Question simple à réponse unique    
Soit $(Z_n)$ une suite géométrique telle que $Z_0=-2$ et de raison $6$, déterminer les variations de $(Z_n)$.
     \begin{reponseshoriz}\bareme{b=1,m=-0.5,e=-0.5,v=0}
      \mauvaise{croissante}
      \bonne{strictement décroissante}
      \mauvaise{strictement croissante}
      \mauvaise{constante}
    \end{reponseshoriz} 
  \end{question}
% b = 1 : une bonne réponse rapporte 1 point
% m = -0.5 : une mauvaise réponse retire 0.5 point
% e = -0.5 : une réponse incohérente retire 0.5 point
% v = 0 : 0 point si aucune réponse  
}


\element{categorie4}{
  \begin{question}{Probleme}    
% Question ouverte   
Le nombre de bactéries présentent dans un organisme suite à une infection est modélisé par $U_{n+1}=U_n\times 1.1$ avec $U_0=10000$ le nombre de bactérie au départ et où $n$ représente le nombre d'heure écoulé.

\begin{enumerate}
\item Calculer $U_1$ et $U_2$.
\item Quelle est la nature de la suite $(U_n)$ ? Précisez sa raison.
\item Calculer le nombre de bactérie présente dans l'organisme au bout de dix heures.
\item Au bout de combien d'heures le nombre de bactérie aura, pour la première fois, doublé ?
\item Au bout de combien d'heures le nombre de bactérie aura, pour la première fois, triplé ?
\end{enumerate}

   \AMCOpen{lines=10,dots=false}{\wrongchoice{NA}\scoring{0}\wrongchoice{AB}\scoring{1}\wrongchoice{B}\scoring{2}\correctchoice{TB}\scoring{3}}
% AB rapporte 1 point
% B rapporte 2 points
% TB rapporte 3 points     
  \end{question}
}

%-------------------------

%%%%%%%%%%%%%%%%%%%%%%%%%
%% fin des questions %%%%
%%%%%%%%%%%%%%%%%%%%%%%%%


%%%%%%%%%%%%%%%%%%%%%%%%%%%%
%%% fabrication des copies%%
%%%%%%%%%%%%%%%%%%%%%%%%%%%%
\exemplaire{18}{    

%%% debut de l'entête des copies :    
	%%% première ligne

	\begin{minipage}{1\linewidth}
	  \centering\large\bf DS 1STMG sur les suites arithmétiques et géométriques 
	\end{minipage}

	%%% codage des numéros d'étudiants sur 2 chiffres par les candidats et encadré d'écriture manuelle des noms.
	\vspace*{3mm}
	{\setlength{\parindent}{0pt}\AMCcode{etu}{2}\hspace*{\fill}
	\begin{minipage}[b]{12cm}
	 \hspace*{5mm}
  	  Codez votre numéro ci contre chiffre par chiffre, puis complétez l'encadré.
     \vspace{3ex}

	 \hfill\champnom{\fbox{
	 \begin{minipage}{1\linewidth}
		\vspace*{3mm}
		NOM - Prénom - Classe :
		\vspace*{3mm}
	 \end{minipage}
	    }
	   }
	\hfill\vspace{1ex}
% texte de présentation 
\begin{center}\em

Durée : 55 minutes.

  Aucun document n'est autorisé. La calculatrice est autorisée.
  Les questions faisant apparaître le symbole \multiSymbole{} peuvent
  présenter une ou plusieurs bonnes réponses. Les autres ont une unique bonne réponse.
  Les réponses fausses ou incohérentes retirent des points.

\end{center}

	\end{minipage}\hspace*{\fill}
	}

\vspace{4ex}

%%% Fin de l'entête
%%%%%%%%%%%%%%%%%%%%%%%

%% Composition des QCM
%%% On mélange des questions par groupe
% On efface le {tout} précédent
\cleargroup{tout}

% On mélange par "catégorie", et on prélève aléatoirement [n] questions 
% du lot {categorie-k} vers la copie complète {tout}
% puis on restitue la copie.
% Si on met [0], toutes les questions sont prises.
\melangegroupe{categorie1}\copygroup[1]{categorie1}{tout}
\melangegroupe{categorie2}\copygroup[1]{categorie2}{tout}
\melangegroupe{categorie3}\copygroup[0]{categorie3}{tout}
\melangegroupe{categorie4}\copygroup[0]{categorie4}{tout}

\restituegroupe{tout}

% Saut d'une page si le sujet comporte un nombre impair de pages.
% Ainsi il est possible d'imprimer les sujets en recto-verso.

\AMCcleardoublepage    

%%% Fin de la création des questionnaires

}
% Fin de exemplaires{k}{ 

\end{document}

